\partieMat{1}{Compter jusqu’à dix}{Dénombrer, constituer, comparer.}{\pic[scale=2.6] at (16.8,6.4) {panier};\pic[scale=1.5,mat/c=matRouge] at (16,7.4) {pomme};\pic[scale=1.5,mat/c=matVert] at (17.5,7.5) {pomme};\pic[scale=1.5,mat/c=matRouge] at (16.7,8.4) {pomme};\pic[scale=1.9,mat/c=matRouge] at (14,9.6) {pomme};\pic[scale=1.4,mat/c=matVert] at (18.8,9.8) {pomme};}{À partir de cinq ans, l’enfant dénombre des collections jusqu’à dix, voire au-delà, en utilisant des compositions (« cinq et deux, sept »). Il constitue des collections à partir de l’écriture chiffrée et compare des nombres.}{\ligneAct{1}{Sept partout}{act-gs-1}\ligneAct{2}{Le bon nombre}{act-gs-2}\ligneAct{3}{Deux groupes de points}{act-gs-3}\ligneAct{4}{Tous les bonbons}{act-gs-4}\ligneAct{5}{Autant que le nombre}{act-gs-5}\ligneAct{6}{Le plus grand nombre}{act-gs-6}\ligneAct{7}{Autant que le modèle}{act-gs-7}}
\begin{activite}{1}{Sept partout}{colorier}{Colorie le rond sous chaque boîte où il y a {\otBold\color{otNavy}sept} objets.}{Reconnaître des collections de sept objets de tailles, de natures et de dispositions différentes.}{act-gs-1}
\matFigurePage{gs-01}
\end{activite}
\begin{activite}{2}{Le bon nombre}{relier}{Compte les objets de chaque boîte. Relie la boîte au {\otBold\color{otNavy}nombre} qui convient.}{Dénombrer une collection jusqu’à dix et l’associer à son écriture chiffrée.}{act-gs-2}
\matFigurePage{gs-02}
\end{activite}
\begin{activite}{3}{Deux groupes de points}{dessiner}{Combien de points sur chaque carte ? {\otBold\color{otNavy}Écris} le nombre dans la case.}{Dénombrer en utilisant des compositions : « quatre et quatre, huit ».}{act-gs-3}
\matFigurePage{gs-03}
\end{activite}
\begin{activite}{4}{Tous les bonbons}{barrer,dire}{Compte les bonbons : fais une petite croix sur chacun quand tu le comptes, pour n’en oublier aucun. Dis combien il y en a.}{Parcourir une grande collection en marquant les éléments déjà comptés.}{act-gs-4}
\matFigurePage{gs-04}
\end{activite}
\begin{activite}{5}{Autant que le nombre}{dessiner}{Lis le nombre. Dessine {\otBold\color{otNavy}autant} de ronds dans la boîte.}{Constituer une collection dont la quantité est donnée par son écriture chiffrée.}{act-gs-5}
\matFigurePage{gs-05}
\end{activite}
\begin{activite}{6}{Le plus grand nombre}{entourer}{Dans chaque cadre, {\otBold\color{otNavy}entoure} le nombre le plus grand. Tu peux t’aider de la bande des nombres.}{Comparer des quantités données par leur écriture chiffrée.}{act-gs-6}
\matFigurePage{gs-06}
\end{activite}
\begin{activite}{7}{Autant que le modèle}{entourer}{Entoure la boîte qui a {\otBold\color{otNavy}autant} de pommes que le modèle.}{Comparer des collections en les dénombrant, quand elles sont organisées différemment.}{act-gs-7}
\matFigurePage{gs-07}
\end{activite}
\partieMat{2}{Composer les nombres}{Un de plus, un de moins, surcompter, les doubles.}{\pic[scale=1.7,mat/n=3] at (14.6,9) {de};\pic[scale=1.7,mat/n=2] at (17.2,7) {de};\pic[scale=2.1,mat/n=3] at (18.6,9.6) {main};\pic[scale=1.3,mat/n=1] at (13.2,6.6) {de};}{L’enfant comprend que chaque nombre s’obtient en ajoutant un au précédent. Il surcompte à partir d’un nombre gardé « dans sa tête » et mémorise des compositions, dont les doubles, qui lui permettent de calculer sans tout recompter.}{\ligneAct{8}{Un de moins}{act-gs-8}\ligneAct{9}{Juste avant, juste après}{act-gs-9}\ligneAct{10}{Je garde dans ma tête}{act-gs-10}\ligneAct{11}{Les doubles des papillons}{act-gs-11}\ligneAct{12}{Les doigts des deux mains}{act-gs-12}\ligneAct{13}{La maison du huit}{act-gs-13}\ligneAct{14}{Les dominos du sept}{act-gs-14}}
\begin{activite}{8}{Un de moins}{barrer,dessiner}{Barre {\otBold\color{otNavy}un} objet dans chaque rangée. Combien en reste-t-il ? Écris le nombre.}{Réaliser une collection qui a un objet de moins et nommer le nombre avant et après le retrait.}{act-gs-8}
\matFigurePage{gs-08}
\end{activite}
\begin{activite}{9}{Juste avant, juste après}{dessiner}{Écris le nombre qui vient {\otBold\color{otNavy}juste avant} et celui qui vient {\otBold\color{otNavy}juste après}.}{Comprendre que chaque nombre s’obtient en ajoutant un au précédent.}{act-gs-9}
\matFigurePage{gs-09}
\end{activite}
\begin{activite}{10}{Je garde dans ma tête}{dessiner,dire}{Garde le nombre du dé dans ta tête, puis compte les jetons en continuant. {\otBold\color{otNavy}Écris} combien il y a en tout.}{Surcompter à partir d’un nombre donné.}{act-gs-10}
\matFigurePage{gs-10}
\end{activite}
\begin{activite}{11}{Les doubles des papillons}{dessiner,dire}{Les deux ailes ont autant de points. Combien de points en tout ? {\otBold\color{otNavy}Écris} le nombre.}{Connaître les doubles : deux et deux, trois et trois, quatre et quatre, cinq et cinq.}{act-gs-11}
\matFigurePage{gs-11}
\end{activite}
\begin{activite}{12}{Les doigts des deux mains}{dessiner}{Combien de doigts sont levés en tout ? {\otBold\color{otNavy}Écris} le nombre.}{Composer les nombres jusqu’à dix avec les doigts des deux mains.}{act-gs-12}
\matFigurePage{gs-12}
\end{activite}
\begin{activite}{13}{La maison du huit}{dessiner}{Dans la maison du huit, chaque étage a huit points en tout. {\otBold\color{otNavy}Dessine} les points qui manquent dans la fenêtre de droite.}{Décomposer huit de plusieurs façons.}{act-gs-13}
\matFigurePage{gs-13}
\end{activite}
\begin{activite}{14}{Les dominos du sept}{entourer}{Entoure les dominos qui ont {\otBold\color{otNavy}sept} points en tout.}{Reconnaître différentes décompositions de sept.}{act-gs-14}
\matFigurePage{gs-14}
\end{activite}
\partieMat{3}{Écrire et réciter}{Les chiffres jusqu’à dix ; la comptine jusqu’à trente.}{\pic[scale=2,mat/c=matJaune] at (14,9.2) {etoile};\pic[scale=1.4,mat/c=matOrange] at (17,7) {etoile};\pic[scale=1.2,mat/c=matJaune] at (18.6,10) {etoile};}{L’enfant écrit les nombres de un à dix. Il récite la comptine jusqu’à trente, à partir d’un nombre donné, de deux en deux et à rebours de dix à un, ce qui prépare le surcomptage et le décomptage.}{\ligneAct{15}{J’écris 7, 8, 9}{act-gs-15}\ligneAct{16}{J’écris jusqu’à 10}{act-gs-16}\ligneAct{17}{Relie de 1 à 20}{act-gs-17}\ligneAct{18}{Le compte à rebours}{act-gs-18}\ligneAct{19}{De deux en deux}{act-gs-19}\ligneAct{20}{Le train part de douze}{act-gs-20}}
\begin{activite}{15}{J’écris 7, 8, 9}{tracer}{Repasse sur les chiffres gris avec ton doigt, puis avec un crayon. Écris le chiffre {\otBold\color{otNavy}tout seul} dans la dernière case.}{Écrire en chiffres les nombres de un à dix.}{act-gs-15}
\matFigurePage{gs-15}
\end{activite}
\begin{activite}{16}{J’écris jusqu’à 10}{tracer,dessiner}{Repasse sur les nombres gris. Puis compte et {\otBold\color{otNavy}écris} le nombre d’objets dans chaque case.}{Écrire en chiffres les nombres de un à dix pour dire une quantité.}{act-gs-16}
\matFigurePage{gs-16}
\end{activite}
\begin{activite}{17}{Relie de 1 à 20}{relier,dire}{Dis la comptine et relie les points dans l’ordre, de 1 à 20. Termine en reliant 20 à 1. Qu’apparaît-il ?}{Connaître la comptine numérique au-delà de dix et suivre la suite des nombres.}{act-gs-17}
\matFigurePage{gs-17}
\end{activite}
\begin{activite}{18}{Le compte à rebours}{dessiner,dire}{La fusée va décoller ! Dis les nombres de dix à un et {\otBold\color{otNavy}écris} ceux qui manquent.}{Réciter la comptine numérique à rebours de dix à un.}{act-gs-18}
\matFigurePage{gs-18}
\end{activite}
\begin{activite}{19}{De deux en deux}{entourer,dire}{Les chaussettes vont par deux. Compte de deux en deux en touchant chaque paire : deux, quatre, six… {\otBold\color{otNavy}Entoure} la paire où tu dis « dix ».}{Réciter la comptine de deux en deux, en partant de deux.}{act-gs-19}
\matFigurePage{gs-19}
\end{activite}
\begin{activite}{20}{Le train part de douze}{colorier,dire}{Le premier wagon porte le nombre douze. Continue la comptine en touchant chaque wagon. Colorie le wagon où tu dis « {\otBold\color{otNavy}seize} ».}{Réciter la comptine en partant d’un nombre autre que un.}{act-gs-20}
\matFigurePage{gs-20}
\end{activite}
\partieMat{4}{Rangs et bande}{Premier, dernier, avant-dernier ; la bande numérique.}{\pic[scale=1.8,mat/c=matRouge] at (14,8) {pion};\pic[scale=1.8,mat/c=matBleu] at (15.8,7.2) {pion};\pic[scale=1.6,mat/n=4] at (18.2,9.2) {de};}{L’enfant repère le rang d’un objet dans une suite de dix, même quand on change le point de départ. Il construit la bande numérique jusqu’à dix et relie l’ajout à un avancement, le retrait à un recul.}{\ligneAct{21}{Dix cartes}{act-gs-21}\ligneAct{22}{Les canards vont à la mare}{act-gs-22}\ligneAct{23}{À la même place}{act-gs-23}\ligneAct{24}{La bande jusqu’à dix}{act-gs-24}\ligneAct{25}{Sur la bande}{act-gs-25}\ligneAct{26}{Avancer, reculer}{act-gs-26}}
\begin{activite}{21}{Dix cartes}{colorier,entourer}{En partant du disque {\otBold\color{otNavy}bleu}, colorie la septième carte. En partant du disque {\otBold\color{otNavy}rouge}, entoure la troisième carte.}{Déterminer un rang dans une suite de dix objets quand on change le point de départ.}{act-gs-21}
\matFigurePage{gs-21}
\end{activite}
\begin{activite}{22}{Les canards vont à la mare}{colorier}{Les canards marchent vers la mare. Colorie le {\otBold\color{otNavy}premier} en rouge, le {\otBold\color{otNavy}deuxième} en jaune, l’{\otBold\color{otNavy}avant-dernier} en vert et le {\otBold\color{otNavy}dernier} en bleu.}{Repérer le premier, le deuxième, l’avant-dernier et le dernier d’une file.}{act-gs-22}
\matFigurePage{gs-22}
\end{activite}
\begin{activite}{23}{À la même place}{dessiner}{Dessine dans le train du bas l’étoile, le rond et le triangle, {\otBold\color{otNavy}dans les mêmes wagons} que sur le modèle.}{Repérer la position d’un objet dans une suite ordonnée pour la reproduire.}{act-gs-23}
\matFigurePage{gs-23}
\end{activite}
\begin{activite}{24}{La bande jusqu’à dix}{dessiner}{Écris les nombres qui {\otBold\color{otNavy}manquent} dans chaque bande.}{Compléter une bande numérique lacunaire jusqu’à dix.}{act-gs-24}
\matFigurePage{gs-24}
\end{activite}
\begin{activite}{25}{Sur la bande}{relier}{Relie chaque carte à la case de la bande qui a le {\otBold\color{otNavy}même nombre}.}{Placer différentes représentations des nombres jusqu’à dix sur la bande numérique.}{act-gs-25}
\matFigurePage{gs-25}
\end{activite}
\begin{activite}{26}{Avancer, reculer}{colorier}{Le pion avance (flèche verte) ou recule (flèche rouge) d’autant de cases que le dé. Colorie la case où il {\otBold\color{otNavy}arrive}.}{Déterminer la position d’arrivée après un déplacement sur la bande numérique.}{act-gs-26}
\matFigurePage{gs-26}
\end{activite}
\partieMat{5}{Petits problèmes}{Réunir, enlever, comparer, regrouper, partager.}{\pic[scale=2.6,mat/c=matBrun] at (15,8) {ours};\pic[scale=2.4,mat/c=matOrange] at (18.2,6.8) {valise};\pic[scale=1.5,mat/c=matJaune] at (18.6,9.8) {ours};}{En fin de maternelle, l’enfant résout des problèmes en une étape en s’aidant d’un dessin, de ses doigts ou d’une décomposition. Certains énoncés sont « discordants » : « de moins » peut demander d’ajouter. On vérifie toujours la réponse.}{\ligneAct{27}{Les cubes dans la boîte}{act-gs-27}\ligneAct{28}{Les oiseaux s’envolent}{act-gs-28}\ligneAct{29}{Une carotte pour chacun}{act-gs-29}\ligneAct{30}{Trois billes de plus}{act-gs-30}\ligneAct{31}{Trois billes de moins}{act-gs-31}\ligneAct{32}{D’où est parti le pion ?}{act-gs-32}\ligneAct{33}{Le goûter}{act-gs-33}\ligneAct{34}{Partager sept gâteaux}{act-gs-34}\ligneAct{35}{Les enveloppes}{act-gs-35}}
\begin{activite}{27}{Les cubes dans la boîte}{dessiner}{On met dans la boîte trois cubes rouges, un cube bleu et deux cubes verts. Combien y a-t-il de cubes dans la boîte ? {\otBold\color{otNavy}Écris} le nombre.}{Déterminer le tout dans un problème à trois parties.}{act-gs-27}
\matFigurePage{gs-27}
\end{activite}
\begin{activite}{28}{Les oiseaux s’envolent}{barrer,dessiner}{Sept oiseaux sont sur la branche. Trois s’envolent. {\otBold\color{otNavy}Barre} les oiseaux qui s’envolent, puis écris combien il en reste.}{Résoudre un problème de retrait en s’aidant d’une représentation.}{act-gs-28}
\matFigurePage{gs-28}
\end{activite}
\begin{activite}{29}{Une carotte pour chacun}{dessiner}{Chaque lapin veut une carotte, mais il n’y en a que deux. {\otBold\color{otNavy}Dessine} les carottes qui manquent, puis écris combien il en faut encore.}{Trouver la quantité à ajouter pour passer d’une collection à une autre.}{act-gs-29}
\matFigurePage{gs-29}
\end{activite}
\begin{activite}{30}{Trois billes de plus}{dessiner}{L’ours a quatre billes. Le lapin a {\otBold\color{otNavy}trois billes de plus} que l’ours. Dessine les billes du lapin.}{Trouver une quantité à partir d’une autre et de l’écart entre les deux.}{act-gs-30}
\matFigurePage{gs-30}
\end{activite}
\begin{activite}{31}{Trois billes de moins}{dessiner}{Le chat a cinq billes. Il a {\otBold\color{otNavy}trois billes de moins} que la souris. Dessine les billes de la souris.}{Résoudre un problème d’écart dont l’énoncé dit « de moins » alors qu’il faut ajouter.}{act-gs-31}
\matFigurePage{gs-31}
\end{activite}
\begin{activite}{32}{D’où est parti le pion ?}{colorier}{Le pion est arrivé sur la case où il est. Il a avancé (flèche verte) ou reculé (flèche rouge) comme le dé. Colorie la case d’où il est {\otBold\color{otNavy}parti}.}{Trouver la position de départ à partir de la position d’arrivée et du déplacement.}{act-gs-32}
\matFigurePage{gs-32}
\end{activite}
\begin{activite}{33}{Le goûter}{dessiner}{On veut mettre {\otBold\color{otNavy}deux} gâteaux dans chaque assiette. Combien faut-il prendre de gâteaux dans la réserve ? Écris le nombre.}{Déterminer le tout dans un problème de groupements.}{act-gs-33}
\matFigurePage{gs-33}
\end{activite}
\begin{activite}{34}{Partager sept gâteaux}{dessiner}{Partage les sept gâteaux entre les deux ours pour qu’ils en aient {\otBold\color{otNavy}autant}. Dessine-les dans les assiettes. S’il en reste un, dessine-le dans la case « le reste ».}{Partager équitablement avec un reste.}{act-gs-34}
\matFigurePage{gs-34}
\end{activite}
\begin{activite}{35}{Les enveloppes}{entourer,dessiner}{On range {\otBold\color{otNavy}deux} images dans chaque enveloppe. Entoure les images deux par deux, puis écris combien il faut d’enveloppes.}{Trouver le nombre de parts dans un problème de groupements.}{act-gs-35}
\matFigurePage{gs-35}
\end{activite}
\partieMat{6}{Formes et solides}{Nommer, décrire, assembler, tracer à la règle.}{\draw[mat contour,fill=matRouge] (15,8) circle (1.2);\draw[mat contour,fill=matJaune] (17.565,8.204) -- (19.196,8.965) -- (18.435,10.596) -- (16.804,9.835) -- cycle;\draw[mat contour,fill=otGreen] (16.815,4.533) -- (19.11,5.138) -- (17.392,6.645) -- cycle;}{L’enfant nomme carré, rectangle, triangle et disque dans toutes les orientations, décrit le cube et la pyramide par leurs faces, reproduit des assemblages (jusqu’à huit pièces) et apprend à tracer à la règle. Aucun solide n’est dessiné en perspective.}{\ligneAct{36}{Quatre familles de formes}{act-gs-36}\ligneAct{37}{Combien de côtés ?}{act-gs-37}\ligneAct{38}{Les faces du cube}{act-gs-38}\ligneAct{39}{Les faces de la pyramide}{act-gs-39}\ligneAct{40}{Le tangram}{act-gs-40}\ligneAct{41}{Avec la règle}{act-gs-41}}
\begin{activite}{36}{Quatre familles de formes}{colorier}{Colorie les carrés en {\otBold\color{otNavy}rouge}, les rectangles en {\otBold\color{otNavy}bleu}, les triangles en {\otBold\color{otNavy}jaune} et les disques en {\otBold\color{otNavy}vert}.}{Reconnaître et classer carrés, rectangles, triangles et disques, quelles que soient leur taille et leur orientation.}{act-gs-36}
\matFigurePage{gs-36}
\end{activite}
\begin{activite}{37}{Combien de côtés ?}{dessiner}{Compte les côtés de chaque forme. {\otBold\color{otNavy}Écris} le nombre dans la case.}{Décrire des formes planes : compter leurs côtés.}{act-gs-37}
\matFigurePage{gs-37}
\end{activite}
\begin{activite}{38}{Les faces du cube}{colorier,manipuler}{On veut fabriquer un cube en collant des faces. {\otBold\color{otNavy}Colorie} toutes les faces qu’il faut. Combien en faut-il ?}{Décrire un cube : six faces carrées, toutes pareilles.}{act-gs-38}
\matFigurePage{gs-38}
\end{activite}
\begin{activite}{39}{Les faces de la pyramide}{colorier,manipuler}{On veut fabriquer une pyramide comme celle de la classe : un carré en dessous et des triangles tout autour. {\otBold\color{otNavy}Colorie} les faces qu’il faut.}{Décrire une pyramide à base carrée : un carré et quatre triangles.}{act-gs-39}
\matFigurePage{gs-39}
\end{activite}
\begin{activite}{40}{Le tangram}{colorier,manipuler}{Le grand carré est fait des mêmes sept pièces que le modèle. {\otBold\color{otNavy}Colorie} chaque pièce de la même couleur que sur le modèle.}{Reproduire un assemblage de sept formes, même s’il n’est pas à la même taille.}{act-gs-40}
\matFigurePage{gs-40}
\end{activite}
\begin{activite}{41}{Avec la règle}{tracer}{Avec ta {\otBold\color{otNavy}règle}, relie les points dans l’ordre : 1 avec 2, 2 avec 3… et le dernier avec 1. Fais la maison, puis l’étoile.}{Utiliser la règle pour tracer des traits droits.}{act-gs-41}
\matFigurePage{gs-41}
\end{activite}
\partieMat{7}{Longueurs et masses}{Bande témoin, rangement, balance.}{\pic[scale=2.4,mat/n=-1] at (16,8.2) {balance};\pic[scale=1.2] at (13.5,9.9) {pasteque};\pic[scale=1.3] at (18.6,10.5) {plume};}{L’enfant compare des longueurs indirectement avec une bande témoin, range jusqu’à cinq objets et reproduit une longueur. Avec la balance de Roberval, il ordonne trois masses et réalise l’équilibre.}{\ligneAct{42}{La bande témoin}{act-gs-42}\ligneAct{43}{Du plus court au plus long}{act-gs-43}\ligneAct{44}{Aussi long}{act-gs-44}\ligneAct{45}{Le plus lourd des trois}{act-gs-45}\ligneAct{46}{L’équilibre}{act-gs-46}}
\begin{activite}{42}{La bande témoin}{manipuler,entourer}{Quel crayon est le plus long ? Marque la longueur de chaque crayon sur une bande de papier pour comparer, puis {\otBold\color{otNavy}entoure} le plus long.}{Comparer des longueurs avec une bande témoin quand on ne peut pas rapprocher les objets.}{act-gs-42}
\matFigurePage{gs-42}
\end{activite}
\begin{activite}{43}{Du plus court au plus long}{dessiner}{Range les rubans du plus court au plus long : écris {\otBold\color{otNavy}1} à côté du plus court, puis 2, 3, 4 et {\otBold\color{otNavy}5} à côté du plus long.}{Ranger cinq objets selon leur longueur.}{act-gs-43}
\matFigurePage{gs-43}
\end{activite}
\begin{activite}{44}{Aussi long}{tracer}{À partir du point, trace avec ta règle un trait {\otBold\color{otNavy}aussi long} que l’objet du dessus.}{Produire un trait de même longueur qu’un objet donné.}{act-gs-44}
\matFigurePage{gs-44}
\end{activite}
\begin{activite}{45}{Le plus lourd des trois}{entourer,barrer}{Regarde les deux balances. En bas, {\otBold\color{otNavy}entoure} le paquet le plus lourd et {\otBold\color{otNavy}barre} le plus léger.}{Ordonner les masses de trois objets à partir de deux comparaisons.}{act-gs-45}
\matFigurePage{gs-45}
\end{activite}
\begin{activite}{46}{L’équilibre}{entourer,manipuler}{Tous les cubes ont la même masse. Combien de cubes faut-il {\otBold\color{otNavy}ajouter} sur le plateau de droite pour que la balance soit droite ? Entoure le nombre.}{Réaliser l’équilibre sur une balance de Roberval.}{act-gs-46}
\matFigurePage{gs-46}
\end{activite}
\partieMat{8}{Les motifs}{Motifs qui grandissent, même règle, invention.}{\pic[scale=1.4,mat/c=matRouge] at (12,7.8) {perle};\pic[scale=1.4,mat/c=matBleu] at (13.1,8.3) {perle};\pic[scale=1.4,mat/c=matRouge] at (14.2,7.8) {perle};\pic[scale=1.4,mat/c=matBleu] at (15.3,8.3) {perle};\pic[scale=1.4,mat/c=matRouge] at (16.4,7.8) {perle};\pic[scale=1.4,mat/c=matBleu] at (17.5,8.3) {perle};\pic[scale=1.4,mat/c=matRouge] at (18.6,7.8) {perle};}{À cinq ans, l’enfant découvre les motifs évolutifs, transcrit un motif avec d’autres symboles, reconnaît des motifs de même structure et en invente. On n’attend pas d’écriture formelle (AABB).}{\ligneAct{47}{Ça grandit}{act-gs-47}\ligneAct{48}{Plus grand, plus petit}{act-gs-48}\ligneAct{49}{Le même motif, autrement}{act-gs-49}\ligneAct{50}{La même règle}{act-gs-50}\ligneAct{51}{J’invente mon collier}{act-gs-51}}
\begin{activite}{47}{Ça grandit}{dessiner,dire}{Dans ces rangées, il y en a {\otBold\color{otNavy}un de plus à chaque fois}. Dessine dans les cases ce qui vient après.}{Repérer et prolonger un motif évolutif.}{act-gs-47}
\matFigurePage{gs-47}
\end{activite}
\begin{activite}{48}{Plus grand, plus petit}{dessiner,dire}{Regarde comment les tours changent. {\otBold\color{otNavy}Dessine} la tour qui vient après dans chaque cadre en pointillés.}{Prolonger un motif évolutif et dire sa règle.}{act-gs-48}
\matFigurePage{gs-48}
\end{activite}
\begin{activite}{49}{Le même motif, autrement}{dessiner}{Refais chaque rangée dans les cases, mais {\otBold\color{otNavy}avec les formes proposées} : la règle doit rester la même.}{Transcrire un motif avec d’autres symboles en gardant sa structure.}{act-gs-49}
\matFigurePage{gs-49}
\end{activite}
\begin{activite}{50}{La même règle}{entourer,dire}{Deux rangées sont faites avec la {\otBold\color{otNavy}même règle}, même si les formes changent. Entoure-les.}{Reconnaître des motifs qui ont la même structure avec des éléments différents.}{act-gs-50}
\matFigurePage{gs-50}
\end{activite}
\begin{activite}{51}{J’invente mon collier}{colorier,dire}{Invente un collier avec une {\otBold\color{otNavy}règle} : choisis tes couleurs et colorie les perles. Puis explique ta règle à un adulte, qui l’écrit.}{Créer un motif et le décrire pour qu’un autre puisse le reproduire.}{act-gs-51}
\matFigurePage{gs-51}
\end{activite}
\partieMat{9}{Se repérer}{Gauche et droite, flèches, quadrillages.}{\pic[scale=2.2] at (14.6,7) {souris};\pic[scale=2] at (18.4,6.6) {fromage};\pic[scale=2] at (16.8,9.6) {arbre};}{L’enfant distingue sa gauche et sa droite, situe les objets les uns par rapport aux autres, code et décode des déplacements avec des flèches et reproduit un dessin sur un quadrillage.}{\ligneAct{52}{Vers la droite}{act-gs-52}\ligneAct{53}{À gauche, à droite}{act-gs-53}\ligneAct{54}{Le code secret}{act-gs-54}\ligneAct{55}{J’écris le chemin}{act-gs-55}\ligneAct{56}{Sur le quadrillage}{act-gs-56}}
\begin{activite}{52}{Vers la droite}{entourer}{Entoure les animaux qui regardent {\otBold\color{otNavy}vers la droite}, comme la flèche.}{Distinguer la droite et la gauche.}{act-gs-52}
\matFigurePage{gs-52}
\end{activite}
\begin{activite}{53}{À gauche, à droite}{entourer,barrer}{Entoure l’objet qui est juste {\otBold\color{otNavy}à droite} du chat. Fais une croix sur l’objet qui est juste {\otBold\color{otNavy}à gauche} du lapin.}{Situer un objet par rapport à un autre avec « à gauche de », « à droite de ».}{act-gs-53}
\matFigurePage{gs-53}
\end{activite}
\begin{activite}{54}{Le code secret}{colorier}{Pars de l’étoile et suis les flèches du code, une case par flèche. {\otBold\color{otNavy}Colorie} la case où tu arrives.}{Réaliser un déplacement à partir de son codage par des flèches.}{act-gs-54}
\matFigurePage{gs-54}
\end{activite}
\begin{activite}{55}{J’écris le chemin}{dessiner}{Le chemin de la souris est en pointillés. {\otBold\color{otNavy}Dessine} une flèche par case dans les cadres du bas pour coder son chemin jusqu’au fromage.}{Coder un déplacement à l’aide de flèches.}{act-gs-55}
\matFigurePage{gs-55}
\end{activite}
\begin{activite}{56}{Sur le quadrillage}{colorier}{Reproduis le dessin du modèle sur le grand quadrillage : {\otBold\color{otNavy}colorie} les mêmes cases, de la même couleur.}{Repérer des cases sur un quadrillage pour reproduire un dessin.}{act-gs-56}
\matFigurePage{gs-56}
\end{activite}
